\documentclass[11pt]{article}
\usepackage[margin=1in]{geometry}
\usepackage{amsmath,amssymb}
\usepackage[hidelinks]{hyperref}

\title{Literature Answer: \(c_0\) Does Not Embed Into \(\mathcal F(\ell_2)\)}
\author{Run \texttt{fa\_banach\_001}}
\date{2026-06-18}

\newcommand{\Free}{\mathcal F}

\begin{document}
\emergencystretch=2em
\maketitle

\paragraph{Status.}
\texttt{literature\_already\_answered}. The extracted open problem from
arXiv:1505.07209 has a negative answer in later literature.

\paragraph{Source question.}
In Cuth--Doucha--Wojtaszczyk, \emph{On the structure of Lipschitz-free spaces},
arXiv:1505.07209, Section ``Open problems'' asks, after Theorem~4.1 on
finite-dimensional Euclidean subsets,
\[
  \text{Is it true that } c_0 \hookrightarrow \Free(\ell_2)?
\]
In the source file this is the unnumbered question at line 349. The later
supporting paper refers to it as Question~3 of Cuth--Doucha--Wojtaszczyk.

\paragraph{Later answer.}
Aliaga--Nous--Petitjean--Prochazka, \emph{Compact reduction in Lipschitz free
spaces}, arXiv:2004.14250, proves compact reduction for weak sequential
completeness. Corollary~2.2 states that if \(X\) is superreflexive, then
\(\Free(X)\) is weakly sequentially complete. Immediately after the corollary
the authors write that, in particular, \(\Free(\ell_2)\) is weakly
sequentially complete and that this provides a negative answer to Question~3
posed by Cuth--Doucha--Wojtaszczyk.

\paragraph{Why this answers the source question.}
Weak sequential completeness passes to closed subspaces. The space \(c_0\) is
not weakly sequentially complete: the sequence of partial sums
\((e_1+\cdots+e_n)_n\) is weakly Cauchy in \(c_0\), since each
\(\ell_1=c_0^*\) functional sees convergent scalar partial sums, but it has no
weak limit in \(c_0\). Therefore a weakly sequentially complete Banach space
cannot contain an isomorphic copy of \(c_0\). Since \(\Free(\ell_2)\) is WSC by
arXiv:2004.14250, one has
\[
  c_0 \not\hookrightarrow \Free(\ell_2).
\]

\paragraph{Scope limitations.}
This is a full negative literature answer for the \(\ell_2\) question from
arXiv:1505.07209. It does not settle the adjacent source question asking
whether \(c_0\hookrightarrow \Free(\ell_1)\).

\paragraph{Search evidence.}
Cheap index search found no prior exact packet for arXiv:1505.07209. Local
corpus searches for \(c_0\), \(\Free(\ell_2)\), weak sequential completeness,
and Lipschitz-free spaces surfaced arXiv:2004.14250, which explicitly names
and answers the source question. Web searches for the exact \(c_0\) and
\(\Free(\ell_2)\) phrases did not surface a newer contradictory status.

\paragraph{References.}
\begin{itemize}
\item M. Cuth, M. Doucha, and P. Wojtaszczyk, \emph{On the structure of Lipschitz-free spaces}, arXiv:1505.07209.
\item R. J. Aliaga, C. Nous, C. Petitjean, and A. Prochazka, \emph{Compact reduction in Lipschitz free spaces}, arXiv:2004.14250.
\end{itemize}

\end{document}
